\subsection{Experimental validation protocol}

To experimentally validate autoreject, our general strategy is to first visually evaluate the
results and thereafter quantify the performance. We describe below the evaluation metric
used, the methods we compare against, and finally the datasets analyzed. All general data
processing was done using the open source software MNE-Python (Gramfort et al., 2013). The code our implementation is based on our developed pyblinker package, which is available at \url{fsdevelopment.github.io/pyblinker}.

\subsection{Evaluation metric}
\label{sec:eval_metric_protocol}

All reported performance numbers use \textbf{event-level} matching between predicted blink
intervals and ground-truth intervals. For each session we compute true positives (TP),
false positives (FP), and false negatives (FN), then derive precision, recall, and F1.
Unless otherwise stated, we report \textbf{macro-averaged} precision/recall/F1 across
sessions (each session contributes one value).

\paragraph{Matching criterion.}
Event matching uses overlap-based IoU thresholding with default IoU = 0.1, consistent with
the stability analysis in Experiment~4 (Section~\ref{sec:exp7}).

\subsection{Competing methods}
\label{sec:competing_methods_protocol}

We compare five conditions throughout the main experiments:
BLINKER-concat, MNE-annot, DBO, Proposed-Mean, and Proposed-Med (primary).
All conditions operate on the same epoch grid (30-second epochs selected by Experiment~1)
and are evaluated with the same event-level scoring.
